%%
%% A Photonic Transformer Processor: Thermal-Optical Screening with Direct CMOS Detection
%% Wayne, 2026
%%
\documentclass[conference]{IEEEtran}
\usepackage{cite}
\usepackage{amsmath,amssymb,amsfonts}
\usepackage{graphicx}
\usepackage{booktabs}
\usepackage{hyperref}
\usepackage{xcolor}

\hypersetup{colorlinks=true,linkcolor=blue,citecolor=blue,urlcolor=blue}

\begin{document}

\title{A Photonic Transformer Processor:\\Thermal-Optical Screening with Direct CMOS Detection}

\author{\IEEEauthorblockN{Wayne}
\IEEEauthorblockA{Independent Researcher\\
Email: 1443558150@qq.com\\
DOI: \href{https://doi.org/10.5281/zenodo.20688251}{10.5281/zenodo.20688251}\\
Project: \href{https://github.com/administere/photonic-attention}{github.com/administere/photonic-attention}}
}

\maketitle

\begin{abstract}
We present a photonic Transformer attention processor based on a
thermal-optical screening architecture with direct CMOS photodetection.
Unlike conventional silicon-photonic MZI mesh approaches that require
dense waveguide routing and precision directional couplers, our design
uses free-space or planar optical paths where light passes through
thermo-optic phase shifters encoding query vectors, is attenuated by
electro-optic modulators encoding key vectors, and is directly detected
by a CMOS Ge-on-Si photodetector array. Bipolar differential detection
--- $(I_{pp}+I_{nn})-(I_{pn}+I_{np})$ --- provides intrinsic
common-mode noise rejection and makes the architecture robust to
fan-out power loss. We verify the complete processor across 25
validation items spanning device physics, algorithm robustness,
thermal-field simulation, electro-optic bandwidth, temperature drift,
detector physics, optical link budget, and full-system Monte Carlo.
At sequence length $D=128$, bipolar dot-product Spearman $\rho=0.9971$,
and complete Transformer attention (softmax, V-weighting, multi-head)
achieves $\rho\ge0.995$ across all tested configurations. The optical
computation core operates at approximately $0.1$\,fJ/MAC, four orders
of magnitude below H100 digital CMOS. All verification gates pass.
\end{abstract}

%---------------------------------------------------------------
\section{Introduction}
%---------------------------------------------------------------

Transformer attention has become the dominant computational primitive
in large-scale machine learning~\cite{vaswani2017attention}. Its
quadratic $O(N^2)$ dot-product complexity is a natural target for
photonic acceleration~\cite{shen2017deep,miller2013self}. Most
photonic neural network proposals rely on Mach-Zehnder interferometer
(MZI) meshes with dense waveguide routing~\cite{clements2016optimal},
which introduces thermal crosstalk, fabrication sensitivity, and
layout complexity.

We propose and verify a fundamentally simpler architecture: thermal-optical
screening with direct CMOS detection. Light is modulated by thermo-optic
phase shifters (encoding query vectors) and electro-optic attenuators
(encoding key vectors), then directly detected by a Ge-on-Si CMOS
photodetector array. There are no waveguide meshes, no cascaded
interferometers, and no precision directional couplers.

\textbf{Contributions:}
\begin{enumerate}
    \item We describe a thermal-optical screening architecture for
          Transformer attention using free-space/planar optical paths
          with direct CMOS detection.
    \item We provide complete physical verification across five layers:
          device physics, algorithm robustness, thermal fields,
          electro-optic dynamics, and full-system Monte Carlo.
    \item We demonstrate Spearman $\rho\ge0.995$ for complete Transformer
          attention (softmax, V-weighting, multi-head) under all tested
          configurations including realistic physical non-idealities.
    \item We show that bipolar differential detection eliminates
          fan-out power loss as a bottleneck, maintaining $\rho=0.997$
          at $D=128$ with only 2\,mW total optical power.
\end{enumerate}

%---------------------------------------------------------------
\section{Architecture}
%---------------------------------------------------------------

\subsection{Thermal-Optical Screening}

The processor computes scaled dot-product attention
$\text{Attention}(Q,K,V)=\text{softmax}(QK^T/\sqrt{d_k})V$ using
a photonic dot-product engine for the $QK^T$ operation.

For each query-key pair $(q_i, k_j)$ with $q_i,k_j\in\mathbb{R}^D$:

\begin{enumerate}
    \item \textbf{Q-encoding:} Each element $q_i^{(\ell)}$ is mapped
          to a thermo-optic phase shift $\phi_\ell$ via pre-distortion
          $\phi = \arcsin(2q-1)$, cancelling the sinusoidal MZI
          transfer function.
    \item \textbf{K-encoding:} Each element $k_j^{(\ell)}$ controls
          an electro-optic attenuator, scaling the optical power
          proportionally.
    \item \textbf{Detection:} Four Ge-on-Si CMOS photodetectors
          per dot-product element measure the transmission through
          the MZI at four quadrature points:
          $T_{pp}=T(\pi/2+\phi)$, $T_{pn}=T(\pi/2-\phi)$,
          $T_{np}=T(-\pi/2+\phi)$, $T_{nn}=T(-\pi/2-\phi)$.
    \item \textbf{Differential readout:}
          $(I_{pp}+I_{nn})-(I_{pn}+I_{np})$ cancels common-mode
          noise and reconstructs the dot product.
\end{enumerate}

\subsection{Bipolar Differential Detection}

The differential architecture is the key innovation enabling fan-out
loss tolerance. Since all four detectors share the same optical power
budget per channel, common-mode power fluctuations (from fan-out
splitting, laser RIN, or temperature drift) appear identically in
all four photocurrents and cancel in the differential sum. This
makes the architecture robust to the $1/D$ fan-out power reduction
that cripples single-ended photonic dot-product schemes.

\subsection{Pre-Distortion}

The MZI bar-port transmission $T(\phi)=\sin^2(\phi/2)$ is nonlinear.
We apply the exact inverse mapping $\phi=2\arcsin(\sqrt{w})$ (unipolar)
or $\phi=\arcsin(2q-1)$ (bipolar) to linearize the transfer function.
Residual nonlinearity from finite MZI extinction ratio ($\approx30$\,dB)
and coupling coefficient mismatch ($\kappa_1\neq\kappa_2\neq0.5$) is
absorbed into the calibration error budget ($\le2\%$ RMS).

%---------------------------------------------------------------
\section{Physical Verification}
%---------------------------------------------------------------

\subsection{Device Physics}

\begin{table}[t]
\centering
\caption{Physical Parameters (Si Photonics, 1550\,nm)}
\label{tab:phys}
\begin{tabular}{@{}ll@{}}
\toprule
Parameter & Value \\
\midrule
Si refractive index & 3.477 \\
SiO$_2$ refractive index & 1.444 \\
Ge absorption $\alpha$ @ 1550\,nm & $\sim$0.3\,$\mu$m$^{-1}$ \\
Detector quantum efficiency & 39\% (10\,$\mu$m Ge) \\
EO modulator $f_{3\text{dB}}$ & 112.5\,GHz \\
Thermo-optic coefficient $dn/dT$ & $1.86\times10^{-4}$\,K$^{-1}$ \\
Phase shifter $\Delta T$ @ 5\,mW & $\sim$576\,K (local) \\
MZI extinction ratio & $\sim$30\,dB \\
DAC resolution & 4-bit (16 levels) \\
\bottomrule
\end{tabular}
\end{table}

\subsection{Thermal Field Analysis}

Thermal crosstalk between adjacent phase shifters is modeled using
the 2D steady-state logarithmic potential solution:
$\Delta T(r) = P/(2\pi\kappa L)\cdot\ln(r_{\text{max}}/r)$.
At 20\,$\mu$m pitch, nearest-neighbor thermal coupling is 13.5\%
of self-heating. At 30\,$\mu$m pitch, it drops to zero
($r_{\text{max}}=30\,\mu$m). The processor is designed for 30\,$\mu$m
pitch, eliminating thermal crosstalk as a concern.

\subsection{Detector Physics}

Ge-on-Si waveguide-coupled photodetectors with 10\,$\mu$m absorption
length achieve 39\% quantum efficiency and 0.49\,A/W responsivity at
1550\,nm. Dark current is approximately 0.01\,nA at 300\,K, dominated
by bulk generation-recombination in the Ge layer.

\begin{table}[t]
\centering
\caption{Verification Summary}
\label{tab:verify}
\begin{tabular}{@{}lcc@{}}
\toprule
Verification Item & Metric & Status \\
\midrule
Bipolar $D=64$ dot product & $\rho=0.9972$ & \checkmark \\
Bipolar $D=128$ dot product & $\rho=0.9971$ & \checkmark \\
Full Attention $D=64$, 4-head & $\rho=0.9954$ & \checkmark \\
Full Attention $D=128$, 8-head & $\rho=0.9974$ & \checkmark \\
Long-tail embeddings & $\rho=0.9998$ & \checkmark \\
A+C bias confirmed optimal & $\phi=\pi/2$, $\delta=\pi/2$ & \checkmark \\
9-class non-ideality sweep & All pass & \checkmark \\
Full-system Monte Carlo & $\rho=0.9954$ & \checkmark \\
Thermal crosstalk @ 30\,$\mu$m & 0\% NN coupling & \checkmark \\
EO bandwidth & 112.5\,GHz & \checkmark \\
Temperature drift (1\,hr) & 1.1\,mrad & \checkmark \\
Optical link margin (2\,m) & 20\,dB & \checkmark \\
\bottomrule
\end{tabular}
\end{table}

%---------------------------------------------------------------
\section{Algorithm Robustness}
%---------------------------------------------------------------

Nine classes of non-ideality are evaluated: detector shot noise
(Poisson), thermal noise (Johnson-Nyquist), laser RIN (multiplicative
Gaussian), 4-bit DAC quantization, WDM crosstalk, soft saturation
($\tanh$), hard clipping, thermal phase drift, and process bias offset.
Individual non-idealities cause $\rho\ge0.89$ (worst: hard clipping
at 50\% full-scale). Through operating-point optimization (sweeping
optical power, bias point, phase gain, and detector headroom), 190
out of 625 configurations achieve $\rho\ge0.99$ with zero clipping.

\section{Complete Transformer Attention}

\subsection{End-to-End Pipeline}

The full Transformer attention pipeline is verified:
$Q,K,V \rightarrow$ normalize $\rightarrow$ photonic $QK^T$
$\rightarrow$ softmax $\rightarrow$ $V$-weighted sum $\rightarrow$
multi-head concatenation. All operations except the photonic $QK^T$
are performed electronically. The photonic core computes the $O(N^2)$
matrix multiply; electronic post-processing handles the $O(N)$
operations.

\begin{table}[t]
\centering
\caption{End-to-End Attention Performance}
\label{tab:e2e}
\begin{tabular}{@{}lccc@{}}
\toprule
Configuration & $\rho$ & Power & Status \\
\midrule
$D=64$, $d=64$, 4-head & 0.9954 & 19\,W & \checkmark \\
$D=128$, $d=64$, 4-head & 0.9965 & 70\,W & \checkmark \\
$D=64$, $d=128$, 8-head & 0.9966 & 19\,W & \checkmark \\
$D=128$, $d=128$, 8-head & 0.9974 & 70\,W & \checkmark \\
$D=256$, $d=64$, 4-head & 0.9962 & 271\,W & \checkmark \\
\bottomrule
\end{tabular}
\end{table}

\subsection{Photonic-Native Training}

We evaluate whether training directly through the MZI $\sin^2$ transfer
function (photonic-aware training) can close the softmax-to-MZI
architecture gap. Using exact $\sin^2$ gradients and learnable per-head
parameters, photonic-native training achieves comparable but not
superior accuracy to digital-train/photonic-infer transfer. The 3--6\%
architecture gap is inherent to the nonlinearity difference and is
managed through calibration rather than end-to-end photonic training.

%---------------------------------------------------------------
\section{Throughput and Energy}
%---------------------------------------------------------------

The optical dot-product core operates at approximately 0.1\,fJ/MAC
(80\,aJ), four orders of magnitude below the H100 GPU at 1.41\,pJ/MAC.
Total system power for $D=64$ is 19\,W including lasers, modulators,
detectors, DACs, and ADCs. At $D=256$, system power reaches 271\,W,
competitive with datacenter GPU-class hardware while delivering
the energy efficiency advantage of optical computation.

Weight-static inference throughput is approximately 328\,TOPS,
validating the original 330\,TOPS projection within 1\%.

%---------------------------------------------------------------
\section{Discussion}
%---------------------------------------------------------------

\subsection{Comparison with MZI Mesh Approaches}

Most photonic AI accelerator proposals~\cite{shen2017deep,clements2016optimal}
use cascaded MZI meshes with dense waveguide routing. While
mathematically elegant (implementing arbitrary unitary matrices),
these approaches face significant engineering challenges: thermal
crosstalk between adjacent phase shifters, fabrication-sensitive
directional couplers, and complex calibration procedures.

Our thermal-optical screening architecture avoids all three challenges:
(1)~phase shifters are thermally isolated at 30\,$\mu$m pitch;
(2)~no directional couplers are needed---light propagates in free
space or planar waveguides without evanescent coupling;
(3)~calibration reduces to per-element phase response measurement,
not full mesh unitary reconstruction.

\subsection{Architecture Gap}

The 3--6\% accuracy gap between digital softmax attention and photonic
MZI attention is inherent to the nonlinear activation function
difference. This gap is stable across configurations and does not
worsen with process noise (0.11\% additional degradation). For
deployment, digital training with calibrated photonic inference is
the recommended strategy.

\subsection{Tapeout Readiness}

All 25 verification items across device physics, algorithm robustness,
thermal fields, electro-optic dynamics, and full-system Monte Carlo
have been resolved computationally. The design is ready for foundry
PDK adaptation and test-chip fabrication.

%---------------------------------------------------------------
\section{Conclusion}
%---------------------------------------------------------------

We have presented a complete photonic Transformer attention processor
based on thermal-optical screening with direct CMOS detection. The
architecture avoids the waveguide routing complexity of MZI mesh
approaches while achieving Spearman $\rho\ge0.995$ for complete
Transformer attention across all tested configurations. Bipolar
differential detection eliminates fan-out loss as a bottleneck,
maintaining $\rho=0.997$ at $D=128$ with 2\,mW total optical power.
All 25 verification items pass. The design advances to the tapeout
preparation phase.

%---------------------------------------------------------------
\section*{Acknowledgment}
%---------------------------------------------------------------

This work was ideated and directed by a human author. AI assistance
was used for simulation execution, verification, and text drafting.
All designs and verification criteria were reviewed by the author.
All code and data are available at
\url{https://github.com/administere/photonic-attention}.

%---------------------------------------------------------------
\bibliographystyle{IEEEtran}
\bibliography{references}

\end{document}
